[1] Lee, Wonkyeong, et al. "Low-dose Computed Tomography Perceptual Image Quality Assessment Grand Challenge." Medical Image Analysis, vol. 99, 2025, p. 103343.
[2] Nirupama, B., et al. "Diagnosis Based Image Quality Assessment and Enhancement for Low Dose CT Image." Frontiers in Radiology, vol. 5, 2025, p. 1704113.
[3] Athar, S., and Z. Wang. "A Comprehensive Performance Evaluation of Image Quality Assessment Algorithms." IEEE Access, vol. 7, 2019, pp. 140030-140070.
[4] Xun, Siyi, et al. "Bibliometric Analysis of CT-IQA." Journal of King Saud University Computer and Information Sciences, vol. 37, no. 92, 2025.
[5] Shen, X., et al. "No-Reference Image Quality Assessment for Low-Dose CT Images Based on Sparse Representation and JND Model." IEEE Access, 2025.
[6] Gao, Q., et al. "CoreDiff: Contextual Error-Modulated Generalized Diffusion Model for Low-Dose CT Denoising and Generalization." IEEE Transactions on Medical Imaging, vol. 43, no. 2, 2024, pp. 745-759.
[7] Zhu, L., et al. "STEDNet: Swin Transformer-Based Encoder–Decoder Network for Noise Reduction in Low-Dose CT." Medical Physics, vol. 50, 2023, pp. 4443-4458.
[8] Luo, L., et al. "HCformer: Hybrid CNN-Transformer Codec Network for Low-Dose CT Denoising." Journal of Data Acquisition and Processing, 2025.
[9] Sori, W. J., et al. "DFD-Net: Lung Cancer Detection from Denoised CT Scan Image Using Deep Learning." Frontiers in Computer Science, vol. 15, 2021, pp. 1-13.
[10] Shan, H., et al. "3-D Convolutional Encoder-Decoder Network for Low-Dose CT via Transfer Learning from a 2-D Trained Network." IEEE Transactions on Medical Imaging, vol. 37, no. 6, 2018, pp. 1522-1534.
[11] Mileto, A., et al. "State of the Art in Abdominal CT: The Limits of Iterative Reconstruction Algorithms." Radiology, vol. 293, no. 3, 2019, pp. 491-503.
[12] Koetzier, L. R., et al. "Deep Learning Image Reconstruction for CT: Technical Principles and Clinical Prospects." Radiology, vol. 306, no. 3, 2023, p. e221257.
[13] McCollough, C. H., et al. "Low Dose CT Image and Projection Data." The Cancer Imaging Archive, 2020.
[14] Willemink, M. J., et al. "Iterative Reconstruction Techniques for Computed Tomography Part 1: Technical Principles." European Radiology, vol. 23, no. 6, 2013, pp. 1623-1631.
[15] Zhang, J., et al. "A Review of Deep Learning Methods for Denoising of Medical Low-Dose CT Images." Computers in Biology and Medicine, vol. 171, 2024, p. 108112.
[16] Chen, H., et al. "Recent Advances and Clinical Applications of Deep Learning in Medical Image Analysis." European Journal of Radiology, vol. 172, 2024, p. 111355.
[17] Wang, Z., and A. C. Bovik. "Modern Image Quality Assessment." Synthesis Lectures on Image, Video, and Multimedia Processing, vol. 2, 2006, pp. 1-156.
[18] Hu, Z., et al. "No-Reference Image Quality Assessment Based on Global Awareness." PLoS ONE, vol. 19, no. 10, 2024, p. e0310206.
[19] Ahmed, N., and S. Asif. "IQA-NRTL: A No-Reference Image Quality Assessment Algorithm Based on Transfer Learning." PeerJ Computer Science, 2024.
[20] Herath, H. M. S. S., et al. "A Systematic Review of Medical Image Quality Assessment." Journal of Imaging, vol. 11, no. 100, 2025.
[21] Hosu, V., et al. "KonIQ-10k: An Ecologically Valid Database for Deep Learning of Blind Image Quality Assessment." IEEE Transactions on Image Processing, vol. 29, 2020, pp. 4041-4056.
[22] Stepien, I., and M. Oszust. "A Brief Survey on No-Reference Image Quality Assessment Methods for Magnetic Resonance Images." Engineering Applications of Artificial Intelligence, vol. 123, 2023, p. 106283.
[23] Dai, Wenliang, et al. "InstructBLIP: Towards General-Purpose Vision-Language Models with Instruction Tuning." arXiv preprint arXiv:2305.06500, 2023.
[24] Radford, A., et al. "Learning Transferable Visual Models From Natural Language Supervision." International Conference on Machine Learning, 2021.
[25] Ke, Junjie, et al. "MUSIQ: Multi-scale Image Quality Transformer." Proceedings of the IEEE/CVF International Conference on Computer Vision, 2021, pp. 5148-5157.
[26] Chen, Z., et al. "IQAGPT: Image Quality Assessment with Vision-Language Models." Visual Computing for Industry, Biomedicine, and Art, vol. 7, no. 20, 2024.
[27] Liu, H., et al. "Visual Instruction Tuning." arXiv preprint arXiv:2304.08485, 2023.
[28] Smits, M., et al. "Minor Head Injury: Guidelines for the Use of CT." Radiology, vol. 245, 2007, pp. 831-838.
[29] Zhang, Q., et al. "Provably Convergent Learned Inexact Descent Algorithm for Low-Dose CT Reconstruction." IEEE Transactions, 2021.
[30] Zhao, Y., et al. "Multi-Task Deep Learning on Medical Imaging." Computers in Biology and Medicine, vol. 153, 2023, p. 106496.
[31] Gao, Q., et al. "Noise-Inspired Diffusion Model for Generalizable Low-Dose CT Reconstruction." Medical Image Analysis, vol. 105, 2025, p. 103710.
[32] Ragoza, M., and K. Batmanghelich. "Physics-Informed Neural Networks." MICCAI 2023, Springer Nature Switzerland AG, 2023, pp. 334-343.
[33] Hounsfield, G. N. "Computerized Transverse Axial Scanning (Tomography): Part 1. Description of System." British Journal of Radiology, vol. 46, 1973, pp. 1016-1022.
[34] P., C., et al. "Observational Studies in Low-Dose CT." MDPI, 2023.
[35] Niu, C., and H. Shan. "Score-Based Generative Models for Medical Imaging." IEEE Transactions, 2023.
[36] Zabic, S., et al. "A Low Dose Simulation Tool for CT Systems with Energy Integrating Detectors." Medical Physics, vol. 40, no. 3, 2013.
[37] Pouget, E., and V. Dedieu. "Comparison of Supervised-Learning Approaches for Designing a Channelized Observer for Image Quality Assessment in CT." Medical Physics, vol. 50, no. 7, 2023, pp. 4282-4295.
[38] Yan, K., et al. "DeepLesion: Automated Mining of Large-Scale Lesion Annotations and Universal Lesion Detection with Deep Learning." Journal of Medical Imaging, vol. 5, no. 3, 2018.
[39] Elhage, N., et al. "A Mathematical Framework for Transformer Circuits." Transformer Circuits Thread, 2021.
[40] Fan, F.-L., et al. "On Interpretability of Artificial Neural Networks: A Survey." IEEE Transactions on Radiation and Plasma Medical Sciences, vol. 5, no. 6, 2021, pp. 741-760.
[41] Li, Y., et al. "MDST: Multi-Domain Sparse-View CT Reconstruction Based on Convolution and Swin Transformer." Physics in Medicine & Biology, vol. 68, 2023, p. 095019.
[42] Greenspan, H., et al., editors. MICCAI 2023 Proceedings Part X. Springer Nature Switzerland AG, 2023.
[43] Goodfellow, I., Y. Bengio, and A. Courville. Deep Learning. MIT Press, 2016.
[44] Wu, T., et al. "A Comprehensive Study of MLLMs for IQA." arXiv preprint, 2024.
[45] Zhou, Z., et al. "Image Quality Evaluation in Deep-Learning-Based CT Noise Reduction Using Virtual Imaging Trial Methods." Medical Physics, vol. 51, 2024, pp. 5399-5413.
[46] Greenspan, H., et al., editors. MICCAI 2023 Proceedings Part II. Springer Nature Switzerland AG, 2023.
[47] Galdran, A., et al. "Multi-Head Multi-Loss Model Calibration." Transactions on Machine Learning Research, 2022.
[48] Sheikh, H. R., et al. "No-Reference Quality Assessment Using Natural Scene Statistics." IEEE Transactions on Image Processing, vol. 14, no. 11, 2005, pp. 1918-1927.
[49] Liu, M., et al. "Blind Image Quality Assessment." IEEE Conference on Computer Vision and Pattern Recognition, 2009.
[50] Cho, H.-H., et al. "Deep Learning Adaptation for Pediatric Cardiac CT Datasets." PLoS ONE, vol. 19, no. 8, 2024, p. e0300090.
